\section{Суперкомпьютерные технологии}

Суперкомпьютерными технологиями называют совокупность аппаратных и программных
средств, предназначенных для решения вычислительно сложных задач с использованием
массового параллелизма.

Типичные области применения суперкомпьютеров:
численное моделирование физических процессов, механика сплошных сред,
климатическое моделирование, молекулярная динамика, обработка больших объёмов
данных, машинное обучение и решение больших задач линейной алгебры.

Главной особенностью высокопроизводительных вычислений является то, что
производительность определяется не только быстродействием отдельного процессора,
но и возможностью эффективно использовать большое число вычислительных устройств,
организовать обмен данными между ними и обеспечить достаточную пропускную
способность памяти.

\subsection{Архитектура высокопроизводительных вычислительных систем}

Современный суперкомпьютер, как правило, представляет собой совокупность
вычислительных узлов, соединённых высокоскоростной сетью.

Типичный вычислительный узел содержит:
\begin{itemize}
	\item один или несколько многоядерных центральных процессоров (CPU);
	\item оперативную память;
	\item локальные уровни кэш-памяти;
	\item сетевой адаптер;
	\item иногда один или несколько графических ускорителей (GPU).
\end{itemize}

Таким образом, современная суперкомпьютерная система обычно является
\textbf{гетерогенной}: разные типы вычислительных устройств предназначены
для различных видов нагрузки.

По характеру организации памяти выделяют две принципиальные архитектуры.

\textbf{Системы с общей памятью} имеют единое адресное пространство:
все процессоры или ядра могут обращаться к одним и тем же данным.

\textbf{Системы с распределённой памятью} состоят из отдельных узлов,
каждый из которых имеет собственную оперативную память. Обмен между узлами
осуществляется посредством передачи сообщений.

На практике распространены гибридные системы:
внутри одного узла имеется общая память, а между узлами память распределена.

Для классификации архитектур также часто используется классификация Флинна:
\begin{itemize}
	\item SISD --- один поток команд, один поток данных;
	\item SIMD --- один поток команд, множество потоков данных;
	\item MISD --- множество потоков команд, один поток данных;
	\item MIMD --- множество потоков команд, множество потоков данных.
\end{itemize}

Современные суперкомпьютеры в целом относятся к классу MIMD, однако внутри
процессоров и ускорителей широко применяется SIMD- или SIMT-параллелизм.

Важнейшей частью суперкомпьютера является сеть межсоединений.
Для неё существенны:
\begin{itemize}
	\item латентность передачи сообщения;
	\item пропускная способность;
	\item топология сети;
	\item способность одновременно обслуживать большое количество обменов.
\end{itemize}

При масштабировании задачи на тысячи вычислительных узлов время коммуникаций
может стать сопоставимым или даже превышать время собственно вычислений.

\subsection{Параллелизм и уровни параллельной обработки}

\textbf{Параллелизм} означает одновременное выполнение нескольких операций.

Можно выделить несколько уровней параллелизма.

\textbf{Параллелизм на уровне инструкций} реализуется непосредственно
процессором: конвейеризация, суперскалярное исполнение, внеочередное выполнение.

\textbf{Векторный параллелизм} состоит в выполнении одной операции сразу
над несколькими элементами данных. Для этого используются SIMD-инструкции.

Например, операция
\[
c_i = a_i + b_i
\]
может выполняться одновременно для нескольких значений $i$.

\textbf{Параллелизм на уровне потоков} реализуется на многоядерных CPU.
Различные потоки программы выполняются различными ядрами.

\textbf{Параллелизм на уровне процессов} используется прежде всего
в системах с распределённой памятью. Каждый процесс имеет собственное адресное
пространство и взаимодействует с другими процессами посредством сообщений.

\textbf{Массовый параллелизм} характерен для GPU, где одновременно могут
выполняться тысячи лёгких вычислительных потоков.

Таким образом, одна программа может одновременно использовать несколько уровней:
например,
\[
\text{MPI} + \text{OpenMP} + \text{SIMD},
\]
или
\[
\text{MPI} + \text{GPU}.
\]

Такой подход называют \textbf{гибридным параллельным программированием}.

\subsection{Иерархия памяти}

Скорость современных процессоров существенно выше скорости доступа
к основной памяти. Поэтому память организуется иерархически.

Типичная иерархия имеет вид:
\[
\text{регистры}
\rightarrow
L1
\rightarrow
L2
\rightarrow
L3
\rightarrow
\text{оперативная память}
\rightarrow
\text{локальное хранилище}
\rightarrow
\text{параллельная файловая система}.
\]

Чем ближе уровень памяти к процессору, тем меньше его объём, но тем ниже
латентность и выше скорость доступа.

Кэш-память автоматически сохраняет недавно использованные участки основной
памяти. Поэтому пространственная и временная локальность обращений к данным
существенно влияет на производительность программы.

Например, последовательная обработка элементов массива обычно значительно
эффективнее случайного доступа к памяти.

В многопроцессорных узлах распространена архитектура NUMA
(\textit{Non-Uniform Memory Access}). Физически память распределена между
процессорными сокетами, хотя логически образует единое адресное пространство.
Доступ ядра к локальной памяти обычно быстрее доступа к памяти другого сокета.

Поэтому для HPC важны:
\begin{itemize}
	\item локальность данных;
	\item привязка потоков к ядрам;
	\item размещение данных в памяти нужного NUMA-узла;
	\item уменьшение количества обращений к памяти;
	\item эффективное использование кэш-памяти.
\end{itemize}

Для оценки вычислительного алгоритма часто используют
\textbf{арифметическую интенсивность}
\[
I =
\frac{\text{число арифметических операций}}
{\text{объём переданных данных}}.
\]

При малой арифметической интенсивности производительность ограничивается
пропускной способностью памяти, а при большой --- производительностью
вычислительных устройств.

В модели Roofline это записывается как
\[
P =
\min\left(P_{\max},\, B_{\max} I\right),
\]
где $P_{\max}$ --- пиковая вычислительная производительность,
$B_{\max}$ --- максимальная пропускная способность памяти.

\subsection{Распределённая память и технология MPI}

В системе с распределённой памятью каждый вычислительный процесс имеет
собственное адресное пространство. Непосредственно обратиться к данным другого
процесса нельзя.

Для взаимодействия используется передача сообщений.

Наиболее распространённым стандартом является
\textbf{MPI} (\textit{Message Passing Interface}).

При запуске MPI-программы создаётся множество процессов.
Каждому процессу присваивается уникальный номер --- \textbf{rank}.

Если используется $p$ процессов, то
\[
\operatorname{rank} \in \{0,1,\ldots,p-1\}.
\]

Процессы объединяются в группы, называемые \textbf{коммуникаторами}.
Основным коммуникатором является \texttt{MPI\_COMM\_WORLD}.

MPI предоставляет два основных типа обменов.

\textbf{Двухточечные операции} передают данные непосредственно между двумя
процессами:
\begin{itemize}
	\item \texttt{MPI\_Send};
	\item \texttt{MPI\_Recv};
	\item неблокирующие \texttt{MPI\_Isend} и \texttt{MPI\_Irecv}.
\end{itemize}

Неблокирующие операции позволяют совмещать вычисления и коммуникации.

\textbf{Коллективные операции} выполняются группой процессов:
\begin{itemize}
	\item \texttt{MPI\_Bcast} --- широковещательная рассылка;
	\item \texttt{MPI\_Scatter} --- распределение данных;
	\item \texttt{MPI\_Gather} --- сбор данных;
	\item \texttt{MPI\_Reduce} --- редукция;
	\item \texttt{MPI\_Allreduce} --- редукция с передачей результата всем процессам;
	\item \texttt{MPI\_Barrier} --- синхронизация.
\end{itemize}

Например, при вычислении суммы
\[
S = \sum_{i=1}^{N} a_i
\]
массив можно разделить между $p$ процессами. Каждый процесс вычисляет локальную
сумму
\[
S_k = \sum_{i\in\Omega_k} a_i,
\]
после чего глобальный результат вычисляется коллективной операцией редукции:
\[
S = \sum_{k=0}^{p-1} S_k.
\]

Основные накладные расходы MPI связаны с передачей сообщений.
Упрощённо время передачи сообщения размера $m$ можно представить как
\[
T_{\mathrm{comm}} = \alpha + \beta m,
\]
где $\alpha$ характеризует латентность сети,
а $\beta$ --- время передачи единицы данных.

Следовательно, большое число маленьких сообщений часто менее эффективно,
чем меньшее число сообщений большего размера.

\subsection{Общая память и технология OpenMP}

При общей памяти все потоки работают с одним адресным пространством.

Одним из основных средств программирования таких систем является
\textbf{OpenMP}.

OpenMP реализует модель \textbf{fork--join}.

Последовательная программа начинает выполняться одним главным потоком.
При входе в параллельную область создаётся группа потоков:
\[
1 \longrightarrow p.
\]
После завершения параллельной области потоки синхронизируются и выполнение
снова продолжается одним потоком:
\[
p \longrightarrow 1.
\]

В языках C и C++ параллельные области задаются директивами вида
\texttt{\#pragma omp ...}.

Основными конструкциями являются:
\begin{itemize}
	\item \texttt{parallel} --- создание параллельной области;
	\item \texttt{for} --- распределение итераций цикла между потоками;
	\item \texttt{sections} --- выполнение разных участков кода разными потоками;
	\item \texttt{reduction} --- параллельное выполнение редукции;
	\item \texttt{critical} --- критическая секция;
	\item \texttt{atomic} --- атомарная операция;
	\item \texttt{barrier} --- барьерная синхронизация.
\end{itemize}

Переменные могут быть:
\begin{itemize}
	\item \texttt{shared} --- общими для потоков;
	\item \texttt{private} --- локальными для каждого потока.
\end{itemize}

Основная проблема программирования с общей памятью ---
\textbf{гонка данных} (\textit{race condition}).

Она возникает, когда несколько потоков одновременно обращаются к одной области
памяти, причём хотя бы один поток изменяет её, а порядок операций не определён.

Использование синхронизации устраняет гонки, но чрезмерная синхронизация
снижает параллельную производительность.

Другой важный эффект --- \textbf{ложное совместное использование}
(\textit{false sharing}). Он возникает, когда разные потоки изменяют разные
переменные, расположенные в одной линии кэша. Несмотря на логическую независимость
данных, между кэшами возникает интенсивный обмен.

\subsection{Вычисления на графических процессорах}

GPU первоначально создавались для компьютерной графики, однако их архитектура
хорошо подходит для массово-параллельных вычислений.

CPU содержит сравнительно небольшое число сложных вычислительных ядер,
оптимизированных для последовательного исполнения и сложного управления.

GPU содержит большое число более простых вычислительных устройств и ориентирован
на выполнение одной и той же последовательности операций над большим количеством
данных.

Для GPU характерна модель SIMT
(\textit{Single Instruction, Multiple Threads}).

Потоки группируются в аппаратные группы --- например, \textit{warp}.
Потоки одной такой группы исполняют одну инструкцию одновременно.

Поэтому условные переходы вида
\[
\texttt{if (condition)}
\]
могут снижать эффективность, если разные потоки одной группы выбирают разные ветви.
Этот эффект называется \textbf{дивергенцией потоков}.

Память GPU также имеет иерархическую структуру:
\begin{itemize}
	\item регистры;
	\item локальная и разделяемая память вычислительного блока;
	\item кэш;
	\item глобальная память GPU.
\end{itemize}

Глобальная память имеет большую пропускную способность, но сравнительно высокую
латентность. Поэтому важным условием эффективных GPU-вычислений является
объединение обращений соседних потоков к соседним участкам памяти.

Распространёнными технологиями программирования GPU являются CUDA,
OpenCL и модели директивного программирования.

Для гибридного суперкомпьютера типична схема
\[
\text{MPI между узлами}
+
\text{GPU внутри узла}.
\]

Особое значение имеет стоимость передачи данных между CPU и GPU.
Если вычислений над данными мало, время копирования может превысить выигрыш
от использования ускорителя.

GPU особенно эффективны для задач с:
\begin{itemize}
	\item большим числом однотипных операций;
	\item высокой степенью параллелизма;
	\item высокой арифметической интенсивностью;
	\item регулярным доступом к памяти.
\end{itemize}

\subsection{Декомпозиция вычислительной задачи}

Для выполнения задачи на параллельной системе её необходимо разбить
на части --- выполнить \textbf{декомпозицию}.

Выделяют несколько основных подходов.

\textbf{Декомпозиция по данным} означает разделение исходных данных между
вычислителями.

Например, область
\[
\Omega
\]
при решении дифференциального уравнения можно представить как
\[
\Omega =
\bigcup_{k=0}^{p-1} \Omega_k.
\]

Каждый процесс обрабатывает собственную подобласть $\Omega_k$.
Значения на границах подобластей периодически передаются соседним процессам.

Такой подход широко применяется в методах конечных разностей,
конечных объёмов и конечных элементов.

\textbf{Функциональная декомпозиция} состоит в том, что различным вычислителям
назначаются разные операции или этапы алгоритма.

\textbf{Задачный параллелизм} используется, если вычисление можно представить
как набор независимых или частично зависимых задач.

К качеству декомпозиции предъявляются несколько требований.

\textbf{Балансировка нагрузки}: вычислительная работа должна распределяться
примерно равномерно:
\[
W_1 \approx W_2 \approx \ldots \approx W_p.
\]

Если один процесс работает значительно дольше остальных, остальные процессы
будут простаивать в ожидании.

\textbf{Минимизация коммуникаций}: следует уменьшать объём и число обменов.

\textbf{Достаточная гранулярность}: вычислительная работа между обменами должна
быть достаточно большой по сравнению со стоимостью коммуникаций.

Условно эффективный параллельный алгоритм должен удовлетворять
\[
T_{\mathrm{calc}} \gg T_{\mathrm{comm}}.
\]

В реальных программах также стремятся совмещать коммуникации и вычисления:
пока часть данных передаётся по сети, процессор выполняет операции над уже
доступными данными.

\subsection{Производительность, ускорение и масштабируемость}

Производительность суперкомпьютеров часто измеряется в FLOPS ---
числе операций с плавающей точкой в секунду.

Используются единицы:
\[
\begin{aligned}
	1~\text{GFLOPS} &= 10^9~\text{FLOPS},\\
	1~\text{TFLOPS} &= 10^{12}~\text{FLOPS},\\
	1~\text{PFLOPS} &= 10^{15}~\text{FLOPS},\\
	1~\text{EFLOPS} &= 10^{18}~\text{FLOPS}.
\end{aligned}
\]

При этом пиковая производительность оборудования и производительность реальной
прикладной программы могут различаться в несколько раз.

Пусть $T_1$ --- время выполнения задачи одним вычислителем,
а $T_p$ --- время выполнения на $p$ вычислителях.

\textbf{Ускорением} называется
\[
S_p = \frac{T_1}{T_p}.
\]

В идеальном случае
\[
S_p = p.
\]

\textbf{Параллельная эффективность} определяется как
\[
E_p = \frac{S_p}{p}
= \frac{T_1}{pT_p}.
\]

Идеальное значение:
\[
E_p = 1.
\]

На практике
\[
E_p < 1
\]
из-за коммуникаций, синхронизаций, дисбаланса нагрузки и последовательных
участков программы.

Пусть доля программы, допускающая параллельное исполнение, равна $\alpha$,
а последовательная доля равна $1-\alpha$.

Тогда согласно \textbf{закону Амдала}
\[
S_p =
\frac{1}
{(1-\alpha)+\dfrac{\alpha}{p}}.
\]

При $p\rightarrow\infty$
\[
S_{\max} =
\frac{1}{1-\alpha}.
\]

Например, если $5\%$ программы остаётся последовательным,
то максимальное ускорение независимо от числа процессоров не превышает
\[
S_{\max} = \frac{1}{0.05}=20.
\]

Закон Амдала относится прежде всего к задаче фиксированного размера.

Однако при увеличении числа процессоров часто увеличивают и размер задачи.
Этот случай описывает \textbf{закон Густафсона}.

Если $s$ --- доля последовательного времени при выполнении задачи на $p$
процессорах, то масштабированное ускорение оценивается как
\[
S_p^{(G)}
= p-s(p-1).
\]

Это показывает, что большое число процессоров эффективно использовать
для решения всё более крупных задач.

Различают два основных типа масштабируемости.

\textbf{Сильная масштабируемость} исследуется при фиксированном размере задачи:
\[
N=\mathrm{const}, \qquad p\uparrow.
\]

Хорошая сильная масштабируемость означает уменьшение $T_p$ при увеличении $p$.

\textbf{Слабая масштабируемость} исследуется при примерно постоянной нагрузке
на один процесс:
\[
\frac{N}{p}\approx\mathrm{const}.
\]

Идеальная слабая масштабируемость означает
\[
T_p \approx \mathrm{const}
\]
при одновременном увеличении $N$ и $p$.

С ростом числа процессов масштабируемость обычно ухудшается из-за:
\begin{itemize}
	\item увеличения числа коммуникаций;
	\item роста стоимости коллективных операций;
	\item синхронизаций;
	\item дисбаланса нагрузки;
	\item недостаточного объёма работы на один процесс.
\end{itemize}

Поэтому использование большего числа процессоров не гарантирует уменьшения
времени расчёта.

\subsection{Отказоустойчивость и воспроизводимость}

При использовании нескольких вычислительных узлов вероятность аппаратного
или программного отказа возрастает с увеличением масштаба системы.

Если расчёт продолжается часы или дни, потеря всей выполненной работы при отказе
одного узла становится неприемлемой.

Основным методом защиты длительных вычислений является
\textbf{checkpoint/restart}.

Через определённые интервалы времени программа сохраняет своё состояние:
\[
X(t_k)\longrightarrow\text{checkpoint}.
\]

После отказа вычисление продолжается не с начального состояния, а с последней
контрольной точки.

Другими методами обеспечения отказоустойчивости являются:
\begin{itemize}
	\item резервирование данных;
	\item репликация вычислений;
	\item перераспределение работы после отказа процесса;
	\item применение отказоустойчивых файловых систем;
	\item алгоритмические методы восстановления данных.
\end{itemize}

Отдельной проблемой является \textbf{воспроизводимость} параллельных вычислений.

Операции с плавающей точкой в общем случае не ассоциативны:
\[
(a+b)+c \neq a+(b+c).
\]

В последовательной программе порядок суммирования обычно фиксирован.
В параллельной редукции порядок операций может зависеть от числа процессов,
реализации библиотеки или порядка прихода сообщений.

Поэтому два математически эквивалентных запуска могут дать немного различные
результаты.

Для повышения воспроизводимости следует:
\begin{itemize}
	\item фиксировать версии компилятора и библиотек;
	\item сохранять параметры компиляции;
	\item фиксировать число процессов и потоков;
	\item фиксировать начальные значения генераторов случайных чисел;
	\item сохранять входные данные и параметры расчёта;
	\item контролировать порядок редукций, если требуется побитовая
	воспроизводимость;
	\item использовать системы контроля версий;
	\item сохранять метаданные вычислительного эксперимента.
\end{itemize}

Таким образом, эффективность суперкомпьютерного расчёта определяется не столько
самим наличием большого числа вычислителей, сколько согласованностью всех
компонентов:
\[
\boxed{
	\text{алгоритм}
	+
	\text{декомпозиция}
	+
	\text{вычислительная архитектура}
	+
	\text{память}
	+
	\text{коммуникации}
}
\]

Хороший параллельный алгоритм должен обеспечивать достаточный параллелизм,
равномерную загрузку вычислителей, локальность данных и относительно малые
накладные расходы на обмены и синхронизацию.

\newpage